% ============================================================
\section{Construção das redes}\label{sec:redes}
% ============================================================

Todas as redes do experimento são grafos de $k$ vizinhos mais próximos (k-NN) sob o cosseno.
Esta seção descreve os conjuntos de vértices, o cálculo das similaridades, os três tratamentos
de empate e as versões direcionada, por união e mútua. A construção é feita pela etapa
\etapa{knn}.

% ------------------------------------------------------------
\subsection{Conjuntos de vértices}\label{sec:red-vertices}

Os tipos da amostra (os \cod{token\_id} distintos entre as ocorrências, cerca de 3,3 mil) saem
só da amostragem, antes de qualquer k-NN. O k-NN roda então sobre três conjuntos de vértices:
\begin{description}
    \item[As $\approx$15 mil ocorrências da amostra.] As redes principais, onde P1--P4 são
    respondidas: lexical por ocorrência (\cod{lex}) e contextuais (\cod{L01}, \cod{L18},
    \cod{L36}), mais as variantes de robustez (\cod{L36n} e as centradas \cod{L01c},
    \cod{L18c}, \cod{L36c}). Elas se restringem à amostra porque só tokens que aparecem num
    texto têm vetor contextual. Todas têm exatamente os mesmos \nVertices{} vértices.
    \item[O vocabulário inteiro.] A rede lexical de tipos (\cod{vocab}), com os 151.669 ids do
    tokenizer menos os 26 especiais, ou seja, 151.643 vértices (\nTiposVocab{} na última
    execução). As 267 linhas de enchimento da matriz de \emph{embeddings} ficam de fora porque
    nenhum texto as produz. Cada vértice tem como atributos a frequência no corpus e na amostra
    (zero para a maioria), a classe de escrita e se aparece no corpus (\nTiposVocabCorpus{}
    tipos).
    \item[Os $\approx$3,3 mil tipos da amostra.] A tabela auxiliar de vizinhos, não reportada
    como rede, usada na derivação exata da rede lexical por ocorrência e no $T_i$ lexical da
    variante (c) (Seções~\ref{sec:vocab-vs-amostra} e~\ref{sec:red-derivacao}).
\end{description}

A classe de escrita de cada tipo do vocabulário é uma de: latina, chinês/japonês/coreano (CJK),
mista (código, misturas de escritas), pontuação, outras escritas (cirílico, árabe\ldots), pedaço
de bytes (decodifica para U+FFFD), espaço, dígitos e especial. Uma medição na fase de
planejamento deu 50\% de tokens de escrita latina, 20\% CJK e 13\% mistos: a maior parte do
vocabulário nunca aparece num texto em português. Tokens pouco treinados podem ter
\emph{embeddings} quase iguais e formar aglomerados ou \emph{hubs} artificiais; isso é
diagnosticado e reportado nos resultados da rede do vocabulário.

% ------------------------------------------------------------
\subsection{Similaridade em \emph{float64}}\label{sec:red-cosseno}

A similaridade entre dois vértices é o cosseno
\begin{equation}
    s_{ij} = \frac{x_i^\top x_j}{\lVert x_i\rVert\,\lVert x_j\rVert}.
\end{equation}
Os vetores são guardados em bf16, convertidos exatamente para \emph{float64}, normalizados
linha a linha e multiplicados em blocos de até 1.024 linhas na CPU (limitados a 512~MiB por
bloco; cerca de 440 linhas na rede do vocabulário). O \emph{float64} é necessário
por causa da tolerância de empate $\eps=10^{-6}$: em \emph{float32}, um produto interno com
$d=2560$ termos acumula erro de arredondamento da ordem de $10^{-6}$ a $10^{-5}$ no cosseno,
comparável à própria tolerância, o que criaria empates falsos e desfaria empates verdadeiros.
Em \emph{float64} o erro fica abaixo de $10^{-12}$.

A mesma rotina serve para a rede do vocabulário: 151,6 mil vetores normalizados ocupam cerca de
3,1 GB em \emph{float64}, e o cálculo em blocos mantém a memória sob controle.

% ------------------------------------------------------------
\subsection{Listas de candidatos}\label{sec:red-candidatos}

Para cada representação, a etapa guarda para cada vértice $i$ uma lista de candidatos ordenada
por similaridade decrescente (e índice crescente), sem o próprio $i$, em formato CSR
(\cod{cand\_<rep>.npz}). A lista tem os $K=128$ mais similares, com uma garantia: ela contém
\emph{todo} $j$ com $s_{ij}\ge s_{i(k_{\max})}-\eps$, onde $s_{i(k_{\max})}$ é a
$k_{\max}$-ésima maior similaridade da linha e $k_{\max}=20$ é o maior $k$ usado. Se um bloco de
empates passar do $K$-ésimo candidato, a linha inteira é recalculada e todos os empatados
entram. Com essa garantia, todos os valores de $k\le k_{\max}$ e todos os tratamentos de empate
saem da mesma lista, sem recalcular similaridades.

% ------------------------------------------------------------
\subsection{Rede direcionada, por união e mútua}\label{sec:red-versoes}

A relação de $k$ vizinhos mais próximos não é simétrica: $j\in\N_i$ não implica $i\in\N_j$. O
objeto produzido pelo k-NN é, portanto, um \textbf{grafo direcionado}
\begin{equation}
    G^{\rightarrow} = \bigl(V,\ \{\,i\to j : j\in\N_i\,\}\bigr),
\end{equation}
em que todo vértice tem grau de saída $k$ (nas variantes de tamanho fixo) e o grau de entrada
varia: vértices escolhidos por muitos outros são \emph{hubs}. A versão direcionada é usada nas
medidas de vizinhança (Jaccard e dominância, que dependem de $\N_i$), na distribuição de graus
de entrada e na reciprocidade.

Para clusterização, distâncias, diâmetro, componentes e comunidades, usamos a versão
\textbf{não direcionada por união},
\begin{equation}
    E^{\cup} = \bigl\{\{i,j\} : j\in\N_i \ \text{ou}\ i\in\N_j\bigr\},
\end{equation}
em que todo vértice tem grau pelo menos $k$ e, portanto, não há vértices isolados (o que não
garante conexidade, como se vê abaixo). A versão
\textbf{mútua},
\begin{equation}
    E^{\cap} = \bigl\{\{i,j\} : j\in\N_i \ \text{e}\ i\in\N_j\bigr\},
\end{equation}
é mais conservadora (só relações confirmadas pelos dois lados, o que corta as arestas que
chegam a \emph{hubs}) e entra como robustez. Distância média e diâmetro são calculados no maior
componente conexo, porque as duas versões podem se fragmentar. A mútua, por construção; a
união, por um mecanismo específico destas redes: um conjunto de mais de $k$ vértices com vetores
idênticos (as ocorrências de um tipo com $f_t>k$ na rede lexical, ou um grupo de prefixo com
mais de $k$ ocorrências nas camadas contextuais) manda todas as suas arestas de saída para
dentro de si e forma um componente fechado se nenhum vértice de fora escolher algum deles. Na
rede lexical isso atinge justamente os tipos frequentes do desenho, e as distâncias do maior
componente os excluem; por isso a tabela de métricas informa, para cada representação, o
número de componentes e a fração de vértices no maior (Seção~\ref{sec:met-estruturais}).

% ------------------------------------------------------------
\subsection{Tratamento de empates}\label{sec:red-empates}

\paragraph{Por que os empates importam.} Na rede lexical por ocorrência, as $f_{t_i}-1$ outras
ocorrências do tipo de $i$ têm $s_{ij}=1$ exatamente. Se $f_{t_i}-1>k$, a escolha de quais delas
entram em $\N_i$ é arbitrária. Se $f_{t_i}-1<k$, as vagas restantes vão para ocorrências do
tipo mais próximo, que também são idênticas entre si e, portanto, também empatam. Empates exatos
também aparecem nas camadas contextuais, nos grupos de prefixo (Seção~\ref{sec:mod-sequencias}).
Como as redes contextuais são comparadas com a lexical, o tratamento de empates afeta
diretamente as medidas de reorganização.

\paragraph{Definição comum.} Seja $s_{i(k)}$ a $k$-ésima maior similaridade da linha $i$. Os
candidatos se dividem em
\begin{equation}
    F_i = \{\,j : s_{ij} > s_{i(k)}+\eps\,\},\qquad
    B_i = \{\,j : |s_{ij}-s_{i(k)}|\le\eps\,\},\qquad
    r_i = k-|F_i|,
\end{equation}
isto é, os que estão com certeza entre os $k$ primeiros ($F_i$, de ``\emph{front}'') e o bloco
empatado na fronteira ($B_i$, de ``\emph{boundary}''). Sempre vale $|F_i|<k\le|F_i|+|B_i|$, e
faltam $r_i$ vizinhos a escolher em $B_i$. Usamos $\eps=10^{-6}$ na configuração principal e
$\eps=10^{-4}$ como teste de sensibilidade. As três variantes são:

\begin{enumerate}[label=(\alph*)]
    \item \textbf{Sorteio com sementes (principal).} $\N_i=F_i\cup R_i$, com $R_i\subset B_i$
    sorteado uniformemente sem reposição, $|R_i|=r_i$. A etapa \etapa{knn} grava o sorteio com
    $R=20$ sementes para toda representação, e as medidas de vizinhança ($J$, $D$, piso de
    ruído) usam todas elas. As métricas estruturais usam as 20 sementes só na rede lexical, onde
    os empates são maciços, e são reportadas como média $\pm$ desvio entre elas; nas camadas
    contextuais os empates vêm só de grupos de prefixo (vetores idênticos), o sorteio quase não
    muda a rede, e as métricas estruturais e as comunidades usam a semente 0, com o desempate
    por posição como checagem. As comunidades da rede lexical usam as sementes 0 e 1 (a segunda
    mostra quanto a partição depende do sorteio). A semente $r$ da
    representação \cod{rep} com $k$ vizinhos usa o gerador
    \cod{default\_rng([438, crc32(rep), k, r])}, de modo que cada combinação tem um fluxo
    aleatório próprio e reprodutível, que não depende da ordem em que as redes são calculadas.
    O desempate pela \textbf{posição} (os $r_i$ menores \cod{occurrence\_id} de $B_i$) é mantido
    como referência determinística; com ele e a versão por união, o k-NN novo reproduz
    exatamente o \cod{build\_union\_knn\_graph} do piloto (teste de regressão).
    \item \textbf{Todos os empatados.} $\N_i=F_i\cup B_i$. Elimina a arbitrariedade, ao custo de
    um grau de saída variável ($|\N_i|\ge k$). Robustez.
    \item \textbf{Tipos distintos.} As ocorrências do próprio tipo são ignoradas, e a vizinhança
    passa a ser um conjunto $T_i$ de $k$ \emph{tipos} $u\neq t_i$. Cada tipo recebe a nota da
    sua ocorrência mais similar a $i$,
    \begin{equation}
        \sigma_{iu} = \max_{j:\,t_j=u} s_{ij},
    \end{equation}
    e $T_i$ são os $k$ tipos de maior nota, com a mesma regra de $F$ e $B$ sobre $\sigma$
    (desempate pelo menor \cod{token\_id} ou por sorteio). Nas camadas contextuais, $T_i$ diz
    \emph{de quais outros tokens} a ocorrência se aproxima no seu contexto. A variante entra só
    nas medidas de vizinhança (Jaccard sobre conjuntos de tipos), porque a dominância lexical
    não faz sentido quando o próprio tipo é excluído.
\end{enumerate}

\paragraph{Dois exemplos na rede lexical.} Com $k=10$:
\begin{itemize}
    \item um tipo com $f_t=50$: $F_i=\emptyset$ e $B_i$ são as 49 outras ocorrências do tipo
    ($s=1$); $r_i=10$. Com o desempate por posição, as 50 ocorrências escolhem as mesmas 10 de
    menor índice, que viram \emph{hubs} com grau de entrada perto de 50, um artefato puro. Com
    o sorteio, cada ocorrência escolhe 10 das 49 ao acaso, e os graus de entrada se espalham em
    torno de 10. Foi esse artefato, visível já no piloto com 104 tokens, que motivou o sorteio
    como regra principal;
    \item um tipo com $f_t=4$ cujo tipo mais próximo $u$ tem $f_u\ge7$: $F_i$ são as 3 outras
    ocorrências do tipo ($s=1$), $B_i$ são as $f_u$ ocorrências de $u$ (todas com a mesma
    similaridade) e $r_i=7$.
\end{itemize}

\tabelaopcional[etapa \etapa{report} a partir de \etapa{knn} (\cod{\_manifest.json})]{knn_empates}
    {Estatísticas de empates por representação e $k$: linhas com empate na fronteira, maior
    bloco empatado, linhas recalculadas por inteiro e tempo de cálculo.}
    {tab:knn-empates}

% ------------------------------------------------------------
\subsection{Derivação exata da rede lexical por ocorrência}\label{sec:red-derivacao}

Na rede lexical, $x_i=e_{t_i}$, então $s_{ij}=c(t_i,t_j)$, onde $c(t,u)$ é o cosseno entre os
\emph{embeddings} dos tipos $t$ e $u$: a similaridade entre duas ocorrências só depende dos
seus tipos. A linha ordenada de $i$ é, portanto:
\begin{enumerate}
    \item as $f_{t_i}-1$ outras ocorrências do mesmo tipo, com $s=1$;
    \item para cada tipo $u$ da amostra, em ordem decrescente de $c(t_i,u)$, as $f_u$
    ocorrências de $u$, todas com a mesma similaridade.
\end{enumerate}
Os candidatos de $i$ (e $\N_i$ em qualquer variante) saem exatamente da matriz completa de
similaridade entre os cerca de 3,3 mil tipos da amostra ($U\times U$, barata nesse tamanho) e
das contagens $f_u$, sem calcular a matriz $15\,000\times15\,000$. Como há um único valor de
similaridade por par de tipos, os blocos de empate são exatos, sem depender de arredondamento.
Um teste compara a derivação com a força bruta. (O parâmetro \texttt{type\_candidates}${}=64$
é outra coisa: o número de tipos candidatos guardados por linha para a variante (c).)

Na variante (c), a nota de um tipo é $\sigma_{iu}=c(t_i,u)$, e $T_i$ é simplesmente o conjunto
dos $k$ tipos da amostra mais próximos de $t_i$ na tabela auxiliar. Quase nunca há empate na
fronteira: empates exatos só entre tipos com \emph{embeddings} idênticos e, com a tolerância
$\eps$, raros casos de tipos distintos com cossenos a menos de $\eps$ um do outro (numa
aproximação com os 3,3 mil tipos mais frequentes da viabilidade, 2 linhas com $k=10$ e 3 com
$k=20$).

% ------------------------------------------------------------
\subsection{Rede lexical de tipos (vocabulário)}\label{sec:red-vocab}

A rede do vocabulário usa os mesmos passos sobre os 151.643 vetores $e_t$: cosseno em
\emph{float64}, candidatos com a garantia de completude e $k\in\{5,10,20\}$, com $k=10$ como
principal. Empates exatos só acontecem entre tipos com \emph{embeddings} idênticos (tokens
nunca treinados, por exemplo): na revisão fixada, 1.903 linhas do vocabulário repetem o vetor
de pelo menos outra linha, em 58 grupos (o maior com 524 tipos). Com a tolerância $\eps$, raros
empates entre tipos distintos também entram em $B$ (cerca de 0,1\% das linhas numa amostra com
$k=10$). Os empates são resolvidos por posição (\cod{token\_id}) e por um sorteio. Um grupo de
vetores idênticos com mais de $k$ tipos manda todas as suas arestas de saída para dentro de si
e pode formar um componente fechado, pelo mesmo mecanismo da
Seção~\ref{sec:red-versoes}. As métricas estruturais usam a versão por união; a
distribuição de graus de entrada, a direcionada.

% ------------------------------------------------------------
\subsection{Configuração principal e grade de robustez}\label{sec:red-robustez}

A configuração principal é fixa: $k=10$, versão por união (direcionada para vizinhanças e graus
de entrada), empates pela variante (a) (20 sementes na rede lexical e nas medidas de
vizinhança, semente 0 nas métricas estruturais das camadas contextuais), $\eps=10^{-6}$ e saídas
brutas dos
blocos. A Tabela~\ref{tab:grade} lista as variações da robustez; os resultados estão na
Seção~\ref{sec:robustez}.

\begin{table}[htbp]
\centering
\caption{Configuração principal e variações de robustez.}\label{tab:grade}
\small
\begin{tabularx}{\linewidth}{@{}l l X@{}}
\toprule
\textbf{Dimensão} & \textbf{Principal} & \textbf{Robustez e motivo} \\
\midrule
$k$ & 10 & 5 e 20: a estrutura não deve depender de um $k$ específico \\
Simetrização & união & mútua: remove arestas que só existem por causa de \emph{hubs} \\
Empates & (a) sorteio, 20 sementes & posição (referência do piloto); (b) todos os empatados;
(c) tipos distintos \\
Tolerância & $\eps=10^{-6}$ & $\eps=10^{-4}$: sensibilidade da definição de empate \\
Bloco 36 & saída bruta (\cod{L36}) & normalizada (\cod{L36n}), o espaço que os
\emph{logits} leem \\
Centragem & não & \cod{L01c}, \cod{L18c}, \cod{L36c}: vetores menos a média sobre as
ocorrências, que remove a direção comum (anisotropia, ativações massivas) \\
Comunidades & Leiden, resolução 1 & resoluções 0,5 e 2; Louvain como checagem \\
Vértices & amostra inteira & só o núcleo ($\approx$11 mil vértices, ainda acima de $10^4$);
planejado: exige um $k$-NN próprio, ainda sem etapa \\
\bottomrule
\end{tabularx}
\end{table}
